\section{Respuestas a las preguntas de investigación}\label{sec:conclusiones:rq}

La Tabla~\ref{tab:conclusiones:rq} responde a las cuatro preguntas de investigación de la Sección~\ref{sec:introduccion:objetivos}.
Cada celda da el resultado principal, y las cifras de detalle están en el capítulo que la sostiene, el que va entre paréntesis detrás de la pregunta.

\begin{table}[tbp]
\centering
\small
\begin{tabular}{|>{\raggedright\arraybackslash}p{3.3cm}|p{10.8cm}|}
\hline
\textbf{Pregunta} & \textbf{Respuesta} \\
\hline
RQ1. Qué APIs de correo incluir y con qué criterios reproducibles (Cap.~\ref{cap:catalogo}) &
El conjunto debe cubrir cuatro formas de cobro: volumen mensual, pago por uso, producto añadido y límite diario.
Sus APIs deben diferir en cómo limitan el envío, por ejemplo con un \textit{rate} por \textit{endpoint} o con un caudal negociado con el proveedor.
Las ocho del catálogo cumplen esa condición.
Es reproducible el protocolo de selección (tres fuentes de inclusión, cinco criterios ordenados y una regla de reemplazo).
La consulta al modelo de lenguaje no lo es, porque repetirla no garantiza las mismas APIs, y queda registrada en el Anexo~\ref{anx:uso-ia}. \\
\hline
RQ2. Si ACM2YAML basta para la capacidad contractual y qué extensiones requiere (Cap.~\ref{cap:evolucion}) &
No basta: las ocho Datasheets dejan en el registro de procedencia, que ninguna herramienta lee, datos como los seis volúmenes contratados de Postmark.
Las extensiones son nueve propuestas.
Siete amplían lo que el esquema expresa, y dos corrigen formas que la guía de ACM2YAML admite y PRIME4API no lee: el límite vacío (P-04) y el precio \texttt{custom} del plan Enterprise de Resend (P-05). \\
\hline
RQ3. Cómo se relacionan los límites de uso y los \textit{rates}, y qué implica para el consumidor (Cap.~\ref{cap:limites}) &
En seis de las ocho APIs la cuota es la restricción activa, el límite que corta antes el envío.
Con los planes que recoge la Datasheet, en lo que el \textit{rate} admite durante la ventana de la cuota cabe entre 103,7 veces la cuota (MailerSend Starter) y 129.600.000 (Postmark Free).
Quien contrata una de esas seis puede dimensionar el volumen sin consultar el \textit{rate} si su cuenta está aprobada y sin restricciones.
En Amazon SES y Mailchimp Transactional no se puede decir que mande la cuota.
El estado de la cuenta cambia la capacidad en seis de las ocho, esas dos incluidas, y ninguna publica el umbral que dispara el cambio. \\
\hline
RQ4. Cómo comparar las APIs con requisitos funcionales y no funcionales a la vez (Cap.~\ref{cap:comparacion}) &
Un procedimiento instancia en cada API el catálogo de 32 preguntas de decisión y responde cada una con el instrumento de su clase: la Datasheet, PRIME4API o HARVEY4API.
Son respondibles por su clase 25, y 21 sin HARVEY4API, que no se ha ejecutado.
El coste, la cuota y el caudal se calculan sobre la Datasheet, y los \textit{webhooks} o la IP dedicada se leen a mano en la matriz de \textit{features}.
En el caso de los 100.000 correos al mes cumplen cuatro APIs, y la más barata es SendGrid Essentials 100K, a 34,95 USD.
Los campos de ACM2YAML v0.3, con un solo volumen por plan, dan 2,4 veces ese precio. \\
\hline
\end{tabular}
\caption[Respuestas a las preguntas de investigación]{Respuesta a cada pregunta de investigación y capítulo que la sostiene.}
\label{tab:conclusiones:rq}
\end{table}
